\section*{Bab \ref{ch:b1:p-block-13-15} --- Blok p I: Golongan Boron, Karbon, dan Nitrogen}

\begin{solution}{exo:b1:p-block-13-15:1}
\ce{NH4NO3}: $-$III dalam \ce{NH4+}, $+$V dalam \ce{NO3-}. \ce{N2O5}: $+$V.
\ce{HNO2}: $+$III. \ce{N2H4}: $-$II. \ce{Mg3N2}: $-$III.
\end{solution}

\begin{solution}{exo:b1:p-block-13-15:2}
\ce{BF3}: boron dengan tiga ikatan dan enam elektron, segitiga datar.
\ce{NH3}: tiga ikatan dan \omterm{def:b1:lewis-resonance-vsepr:lewis}{sepasang elektron bebas}. \omterm{def:b1:lewis-resonance-vsepr:lewis}{Pasangan elektron bebas}
N membentuk \omterm{def:b1:lewis-resonance-vsepr:lewis-acid}{ikatan datif} dengan B: \ce{BF3} adalah \omterm{def:b1:lewis-resonance-vsepr:lewis-acid}{asam Lewis}, \ce{NH3}
\omterm{def:b1:lewis-resonance-vsepr:lewis-acid}{basa Lewis}; dalam senyawa adisi itu kedua atom tetrahedral.
\end{solution}

\begin{solution}{exo:b1:p-block-13-15:3}
\ce{CO2 + 2NaOH -> Na2CO3 + H2O}; \ce{SiO2 + 2NaOH -> Na2SiO3 + H2O};
\ce{P4O10 + 12NaOH -> 4Na3PO4 + 6H2O}; \ce{Al2O3 + 6HCl -> 2AlCl3 + 3H2O};
\ce{Al2O3 + 2NaOH + 3H2O -> 2NaAl(OH)4}.
\end{solution}

\begin{solution}{exo:b1:p-block-13-15:4}
$\Delta_r G^\circ = 2 \times (-16.45) = \qty{-32.9}{kJ/mol}$, $\log K =
32.9/5.708 = 5.76$, $K = 10^{5.76}$.
\end{solution}

\begin{solution}{exo:b1:p-block-13-15:5}
Cincin enam satuan \ce{SiO3^2-}: \ce{Si6O18^12-}. Rantai ganda, per empat
silikon: dua dengan dua sudut bersama (3 O, muatan masing-masing $-2$),
dua dengan tiga ($2.5$ O, muatan masing-masing $-1$): \ce{Si4O11^6-}.
\end{solution}

\begin{solution}{exo:b1:p-block-13-15:6}
\qty{550}{kg} \ce{Al(OH)3} adalah $550/78.0 = \qty{7.05}{kmol}$, yang
menghasilkan \qty{3.53}{kmol} \ce{Al2O3}, \qty{360}{kg}; pada
\qty{92}{\percent}: \qty{331}{kg}.
\end{solution}

\begin{solution}{exo:b1:p-block-13-15:7}
Boron mempunyai orbital kosong dan menerima \omterm{def:b1:lewis-resonance-vsepr:lewis}{pasangan elektron bebas} ion
hidroksida yang diambil dari air: \ce{B(OH)3 + 2H2O <=> B(OH)4- + H3O+}.
Proton yang dilepaskan berasal dari molekul air, bukan dari asam borat.
\end{solution}

\begin{solution}{exo:b1:p-block-13-15:8}
N berubah dari $-$III menjadi $+$II, lima elektron per \ce{NH3}; setiap
\ce{O2} mengambil empat. $4 \times 5 = 5 \times 4$: \ce{4NH3 + 5O2 -> 4NO + 6H2O}.
\end{solution}

\begin{solution}{exo:b1:p-block-13-15:9}
\ce{NH3 + HNO3 -> NH4NO3}. Satu ton adalah $10^6/80.0 = \qty{1.25e4}{mol}$;
\qty{1.25e4}{mol} amonia untuk penetralan dan \qty{1.25e4}{mol} untuk
membuat asam nitratnya: $2.50 \times 10^4 \times 17.0 = \qty{425}{kg}$.
\end{solution}

\begin{solution}{exo:b1:p-block-13-15:10}
Atom-atom \omterm{def:b1:periodicity:table}{periode} kedua berukuran kecil: orbital \textit{p}-nya tumpang
tindih dengan baik secara berdampingan dan membentuk ikatan $\pi$ yang
kuat, sehingga \ce{N#N} dan \ce{O=C=O} adalah molekul kecil yang stabil.
Atom-atom \omterm{def:b1:periodicity:table}{periode} ketiga lebih besar; tumpang tindih $\pi$-nya buruk, dan
atom-atom itu lebih menyukai lebih banyak ikatan tunggal: tetrahedron
\ce{P4}, jaringan \ce{SiO4}.
\end{solution}

\begin{solution}{exo:b1:p-block-13-15:11}
$E^\circ(\ce{PbO2}/\ce{Pb^2+}) = 1.46 > E^\circ(\ce{Cl2}/\ce{Cl-}) = 1.36$~V:
\ce{PbO2 + 4H+ + 2Cl- -> Pb^2+ + Cl2 + 2H2O}. Timbal(IV) tidak stabil
terhadap timbal(II) (pasangan inert); silikon(II) bukan keadaan yang
stabil dan \ce{SiO2} bukan \omterm{def:b1:nernst:couple}{oksidator}.
\end{solution}

\begin{solution}{exo:b1:p-block-13-15:12}
$160 \times 17.0/14.0 = \qty{194}{Mt}$ amonia, yang mengandung $194 \times
3.0/17.0 = \qty{34}{Mt}$ hidrogen, yang sebagian besar dibuat dari gas alam
(metana dan uap air).
\end{solution}

\begin{solution}{pb:b1:p-block-13-15:1}
\textbf{1.} :N\,$\equiv$\,N: --- ikatan rangkap tiga, sangat kuat, tanpa
kepolaran.
\textbf{2.} 0, $-$III, $+$II, $+$IV, $+$V, serta $-$III/$+$V dalam \ce{NH4NO3}.
\textbf{3.} Bakteri pada akar kacang-kacangan, petir, proses Haber--Bosch.
\textbf{4.} Pemutusan ikatan rangkap tiga memerlukan \omterm{def:b1:elementary-steps:catalyst}{katalis} yang tidak
dimiliki tumbuhan; bakteri yang mempunyai enzim nitrogenase memilikinya.
\textbf{5.} Nitrogen direduksi (0 menjadi $-$III), hidrogen dioksidasi (0
menjadi $+$I).
\textbf{6.} \ce{N2 + 3H2 <=> 2NH3}.
\textbf{7.} $K = 10^{2 \times 16.45/5.708} = 10^{5.76}$.
\textbf{8.} Pada \qty{25}{\celsius} reaksinya jauh terlalu lambat;
\omterm{def:b1:elementary-steps:catalyst}{katalisnya} hanya bekerja dalam keadaan panas, dan kondisinya merupakan
kompromi yang dibahas dalam jilid Tahun~2.
\textbf{9.} \qty{58.8}{kmol} amonia: \qty{29.41}{kmol} \ce{N2},
\qty{824}{kg}; \qty{88.2}{kmol} \ce{H2}, \qty{176}{kg}.
\textbf{10.} $V(\ce{N2}) = 29.4 \times 10^3 \times 8.314 \times 298.15/10^5 =
\qty{729}{m^3}$; udara $729/0.78 = \qty{935}{m^3}$.
\textbf{11.} Dari gas alam (reformasi kukus metana), dengan produk samping
\ce{CO2}.
\textbf{12.} Jika diinjeksikan ke dalam tanah, amonia adalah pupuk yang
murah, tetapi amonia adalah gas beracun; padatan (urea, amonium nitrat)
lebih mudah disimpan dan ditebarkan.
\textbf{13.} N: $-$III $\to$ $+$II (5 elektron), O: 0 $\to$ $-$II (4 per
\ce{O2}): \ce{4NH3 + 5O2 -> 4NO + 6H2O}.
\textbf{14.} \ce{2NO + O2 -> 2NO2}.
\textbf{15.} \ce{3NO2 + H2O -> 2HNO3 + NO}: sebuah \omterm{def:b1:e-ph-diagrams:disproportionation}{disproporsionasi} (N $+$IV
menjadi $+$V dan $+$II).
\textbf{16.} \ce{NH3 + 2O2 -> HNO3 + H2O}.
\textbf{17.} $2 \times 58.8 \times 32.0 = \qty{3.76}{t}$.
\textbf{18.} Kasa itu membuat oksidasi berlangsung cepat dan selektif
menuju NO, bukan menuju \ce{N2} yang lebih stabil.
\textbf{19.} \ce{NH3 + HNO3 -> NH4NO3}.
\textbf{20.} Separuhnya.
\textbf{21.} \qty{29.4}{kmol} \ce{NH4NO3}, \qty{2.35}{t}.
\textbf{22.} $28.0/80.0 = \qty{35}{\percent}$.
\textbf{23.} Zat itu mengandung \omterm{def:b1:nernst:couple}{oksidator} (nitrat) dan \omterm{def:b1:nernst:couple}{reduktor} (amonium)
dalam \omterm{def:b1:crystals-metals:crystal}{kristal} yang sama: jika dipanaskan atau dicampur dengan bahan
bakar, zat itu dapat terurai secara eksplosif.
\textbf{24.} $58.8 \times 63.0 =$ \textbf{\qty{3.7}{t} asam nitrat per ton
amonia}.
\end{solution}
